% Generated draft: requires mathematical and source review.
\documentclass[11pt]{article}
\usepackage[margin=25mm]{geometry}
\usepackage{fontspec}
\setmainfont{Latin Modern Roman}
\usepackage{amsmath,amssymb,mathtools,unicode-math}
\setmathfont{Latin Modern Math}
\usepackage{enumitem,needspace,xurl,hyperref}
\setlength{\parindent}{0pt}
\setlength{\parskip}{5pt}
\setlength{\emergencystretch}{3em}
\newcommand{\N}{\mathbb{N}}
\newcommand{\Z}{\mathbb{Z}}
\newcommand{\Q}{\mathbb{Q}}
\newcommand{\R}{\mathbb{R}}
\newcommand{\C}{\mathbb{C}}
\newcommand{\F}{\mathbb{F}}
\newcommand{\A}{\mathbb{A}}
\newcommand{\PP}{\mathbb{P}}
\newcommand{\E}{\mathbb{E}}
\newcommand{\eps}{\varepsilon}
\newcommand{\dd}{\,\mathrm{d}}
\begin{document}
\section*{UnsolvedMath ID: 3700037}
\textbf{Source problem code:} catalogue entry
\subsection{Definitions and mathematical statement}
For an integer $d\geq1$, let $p(z)=z^d+a_{d-1}z^{d-1}+\cdots+a_0$ be monic with complex coefficients. Its unit lemniscate is
\[
 E(p)=\{z\in\C:|p(z)|=1\}.
\]
Let $\ell(E(p))=\mathcal H^1(E(p))$ be its geometric arclength, including all components and allowing singular level points, and put
\[
 L_d=\sup\{\ell(E(p)):p\in\C[z]\text{ monic},\ \deg p=d\}.
\]
\textbf{Problems.} Determine the maximal length $L_d$, including the meaning of maximum/attainment if required. Is the candidate $p_d(z)=z^d+1$ asymptotically extremal, with the source's stated maximal-length asymptotic
\[
 L_d=2d+o(1)\qquad(d\to\infty)?
\]
The displayed little-$o$ is transcribed literally: it asserts the additive difference $L_d-2d\to0$, which is stronger than $L_d/(2d)\to1$. Likewise, ``asymptotically extremal'' could mean a relative ratio tending to one or an additive extremal gap tending to zero; the excerpt does not distinguish these. The precise error term for the candidate and the maximum must be checked rather than silently replaced by a different asymptotic.

The supplied research summary labels the Erd\H{o}s--Herzog--Piranian conjecture solved for all sufficiently large degrees, attributing this to Tao (2025), and claims confirmation of the printed asymptotic. It also says a bounded, feasibly decidable set of small degrees remains to be checked using a reduction in that work. Consequently the summary's solved label does not itself provide exact values or verified completion in every degree. No theorem statement, error term, small-degree bound, or full citation is supplied here.

The background says the record was retained from Eremenko's author-maintained unsolved index after applying explicit solved/counterexample updates through its reported access date. That provenance and the later solved claim are retained separately, without independently assessing the cited result.
\subsection{Short English statement}
Maximize the arclength of the unit lemniscate of a monic degree-$d$ polynomial and test the proposed extremal shape $z^d+1$. The source prints an additive $2d+o(1)$ asymptotic; its exact error term and the claimed large-degree solution need checking, especially because the summary leaves small degrees unresolved.
\subsection{Sources}
\begin{itemize}
\item Local UnsolvedMath registry, permanent ID 3700037, source code catalogue entry, ``Maximum length of a polynomial lemniscate''; attribution field: Alexandre Eremenko and cited proposers; supplied proposed year 2015. Status solved, classification SOLVED-IN-LITERATURE, updated 2026-08-06.
\item Eremenko, \emph{My favorite unsolved problems}, item ``Erd\H{o}s lemniscate problem,'' \url{https://www.math.purdue.edu/~eremenko/uns1.html}. The record reports PDF-text extraction and access on 2026-07-29, with explicit index updates applied. Rights: NEEDS\_REVIEW; public access verified but redistribution permission not established.
\item The research summary names the Erd\H{o}s--Herzog--Piranian conjecture and attributes a sufficiently-large-degree result and small-degree reduction to Tao (2025), without title, URL, theorem number, or explicit degree bound.
\end{itemize}
\end{document}
